\subsection{Orthogonal Isometries and Symmetry Groups}

Throughout, all spaces are finite-dimensional over a fixed field $\mathbb{F}$; characteristic hypotheses are stated locally.

\begin{proposition} \label{prop:orthogonal-isometry-preserves-quadratic-form} \hlt{Isometries from Quadratic Values}\\
Let $\mathcal{V}$ be nonsingular orthogonal with $\operatorname{char}(\mathbb{F})\neq2$, and put $q(\mathbf{x})\defeq\innerprod{\mathbf{x},\mathbf{x}}$. For $T\in\operatorname{GL}(\mathcal{V})$, $T\in\mathcal{O}(\mathcal{V})\Longleftrightarrow(\forall\,\mathbf{x}\in\mathcal{V})\,q(T\mathbf{x})=q(\mathbf{x})$.
\begin{proof}
\begin{enumerate}[label=\roman*.]
\setlength{\itemsep}{0pt}
\item[$(\Rightarrow)$] If $T\in\mathcal{O}(\mathcal{V})$, then $(\forall\,\mathbf{x}\in\mathcal{V})\,q(T\mathbf{x})=\innerprod{T\mathbf{x},T\mathbf{x}}=\innerprod{\mathbf{x},\mathbf{x}}=q(\mathbf{x})$.
\item[$(\Leftarrow)$] Assume $q\circ T=q$. For $\mathbf{x},\mathbf{y}\in\mathcal{V}$, $2\innerprod{T\mathbf{x},T\mathbf{y}}=q(T(\mathbf{x}+\mathbf{y}))-q(T\mathbf{x})-q(T\mathbf{y})$ by linearity and polarization.\\
Thus $2\innerprod{T\mathbf{x},T\mathbf{y}}=q(\mathbf{x}+\mathbf{y})-q(\mathbf{x})-q(\mathbf{y})=2\innerprod{\mathbf{x},\mathbf{y}}$ by $q\circ T=q$ and polarization.\\
Hence $\innerprod{T\mathbf{x},T\mathbf{y}}=\innerprod{\mathbf{x},\mathbf{y}}$ by $\operatorname{char}(\mathbb{F})\neq2$, and $T\in\mathcal{O}(\mathcal{V})$ by $T\in\operatorname{GL}(\mathcal{V})$.
\end{enumerate}
\hlt{Structural Commentary}: Diagonal preservation and polarization transfer quadratic invariance to every cross-pairing.\\
\hlt{Key Insight}: Invertibility of $2$ makes the diagonal determine the symmetric form.
\end{proof}
\end{proposition}

\begin{definition} \label{def:orthogonal-symmetry-rotation-reflection} \hlt{Special Orthogonal Group and Symmetries}\\
Let $\mathcal{V}$ be nonsingular orthogonal with $\operatorname{char}(\mathbb{F})\neq2$, and put $q(\mathbf{x})\defeq\innerprod{\mathbf{x},\mathbf{x}}$.
\begin{enumerate}[label=(\roman*)]
\setlength{\itemsep}{0pt}
\item The \hlt{special orthogonal group} is $\operatorname{SO}(\mathcal{V})\defeq\{\sigma\in\mathcal{O}(\mathcal{V}):\det(\sigma)=1\}$; its elements are \hlt{rotations}.
\item If $q(\mathbf{u})\neq0$, the \hlt{symmetry} with normal line $\mathbb{F}\mathbf{u}$ is the map $\sigma_{\mathbf{u}}:\mathcal{V}\to\mathcal{V}$. It is defined by $(\forall\,\mathbf{x}\in\mathcal{V})\,\sigma_{\mathbf{u}}(\mathbf{x})\defeq\mathbf{x}-2\frac{\innerprod{\mathbf{x},\mathbf{u}}}{q(\mathbf{u})}\mathbf{u}$.
\end{enumerate}
\end{definition}

\begin{proposition} \label{prop:symmetry-matches-equal-length-nonisotropic-vectors} \hlt{Symmetries and Equal Quadratic Values}\\
Let $\mathcal{V}$ be nonsingular orthogonal with $\operatorname{char}(\mathbb{F})\neq2$, and put $q(\mathbf{x})\defeq\innerprod{\mathbf{x},\mathbf{x}}$.
\begin{enumerate}[label=(\roman*)]
\setlength{\itemsep}{0pt}
\item If $q(\mathbf{w})\neq0$, then $\sigma_{\mathbf{w}}\mathbf{w}=-\mathbf{w}$ and $\sigma_{\mathbf{w}}|_{(\mathbb{F}\mathbf{w})^{\perp}}=I$.
\item If $q(\mathbf{w})\neq0$, then $\sigma_{\mathbf{w}}\in\mathcal{O}(\mathcal{V})$ and $\sigma_{\mathbf{w}}^2=I$.
\item If $q(\mathbf{u})=q(\mathbf{v})\neq0$, then $q(\mathbf{u}-\mathbf{v})\neq0$ or $q(\mathbf{u}+\mathbf{v})\neq0$. In the respective cases, $\sigma_{\mathbf{u}-\mathbf{v}}\mathbf{u}=\mathbf{v}$ and $\sigma_{\mathbf{u}+\mathbf{v}}\mathbf{u}=-\mathbf{v}$.
\end{enumerate}
\begin{proof}
\begin{enumerate}[label=(\roman*)]
\setlength{\itemsep}{0pt}
\item If $q(\mathbf{w})\neq0$, then $\sigma_{\mathbf{w}}\mathbf{w}=\mathbf{w}-2\mathbf{w}=-\mathbf{w}$. For $\mathbf{z}\in(\mathbb{F}\mathbf{w})^{\perp}$, $\innerprod{\mathbf{z},\mathbf{w}}=0$, so $\sigma_{\mathbf{w}}\mathbf{z}=\mathbf{z}$.
\item The formula for $\sigma_{\mathbf{w}}$ is linear, and $\mathcal{V}=\mathbb{F}\mathbf{w}\odot(\mathbb{F}\mathbf{w})^{\perp}$ by Theorem \ref{thm:metric-orthogonal-complement-dimension}(ii).\\
Thus $\sigma_{\mathbf{w}}(\alpha\mathbf{w}+\mathbf{z})=-\alpha\mathbf{w}+\mathbf{z}$ by (i), so $\sigma_{\mathbf{w}}^2=I$ and $q\circ\sigma_{\mathbf{w}}=q$. Hence $\sigma_{\mathbf{w}}\in\mathcal{O}(\mathcal{V})$ by Proposition \ref{prop:orthogonal-isometry-preserves-quadratic-form}.
\item Put $\mathbf{a}=\mathbf{u}-\mathbf{v}$ and $\mathbf{b}=\mathbf{u}+\mathbf{v}$. Then $\innerprod{\mathbf{a},\mathbf{b}}=q(\mathbf{u})-q(\mathbf{v})=0$.\\
The conjunction $q(\mathbf{a})=q(\mathbf{b})=0$ would imply $0=q(\mathbf{a}+\mathbf{b})=q(2\mathbf{u})=4q(\mathbf{u})$, a contradiction.\\
Therefore $q(\mathbf{a})\neq0$ or $q(\mathbf{b})\neq0$. If $q(\mathbf{a})\neq0$, then $q(\mathbf{a})=2\innerprod{\mathbf{u},\mathbf{a}}$ and $\sigma_{\mathbf{a}}\mathbf{u}=\mathbf{u}-\mathbf{a}=\mathbf{v}$.\\
If $q(\mathbf{b})\neq0$, then $q(\mathbf{b})=2\innerprod{\mathbf{u},\mathbf{b}}$ and $\sigma_{\mathbf{b}}\mathbf{u}=\mathbf{u}-\mathbf{b}=-\mathbf{v}$.
\end{enumerate}
\hlt{Structural Commentary}: The normal-line decomposition verifies a symmetry; the sum-difference pair then aligns equal quadratic values.\\
\hlt{Key Insight}: Orthogonal sum-difference directions cannot both be isotropic.
\end{proof}
\end{proposition}

\begin{theorem} \label{thm:orthogonal-transformations-products-of-symmetries} \hlt{Orthogonal Transformations Are Products of Symmetries}\\
Let $\mathcal{V}$ be nonsingular orthogonal with $\operatorname{char}(\mathbb{F})\neq2$, and put $q(\mathbf{x})\defeq\innerprod{\mathbf{x},\mathbf{x}}$. Then $\mathcal{O}(\mathcal{V})=\langle\sigma_{\mathbf{u}}:\mathbf{u}\in\mathcal{V},\ q(\mathbf{u})\neq0\rangle$.
\begin{proof}
\begin{enumerate}[label=(\roman*)]
\setlength{\itemsep}{0pt}
\item \hlt{Induction Basis}: Induct on $n=\dim\mathcal{V}$. If $n=0$, then $\mathcal{O}(\mathcal{V})=\{I\}$, generated by the empty product.\\
If $n=1$, fix $\mathbf{0}\neq\mathbf{u}\in\mathcal{V}$. For $\tau\in\mathcal{O}(\mathcal{V})$, write $\tau\mathbf{u}=\lambda\mathbf{u}$.\\
Nonsingularity implies $q(\mathbf{u})\neq0$, and $\lambda^2=1$ by $q(\tau\mathbf{u})=q(\mathbf{u})$. Thus $\tau=I$ or $\tau=\sigma_{\mathbf{u}}$ by Proposition \ref{prop:symmetry-matches-equal-length-nonisotropic-vectors}(i).
\item \hlt{Nonisotropic Line}: Assume $n>1$. If $(\forall\,\mathbf{x}\in\mathcal{V})\,q(\mathbf{x})=0$, then the form is alternating.\\
It is zero by Proposition \ref{prop:form-type-relations-by-characteristic}(ii), contrary to nonsingularity. Thus $(\exists\,\mathbf{u}\in\mathcal{V})\,q(\mathbf{u})\neq0$; fix such a $\mathbf{u}$.
\item \hlt{Line Alignment}: Fix $\tau\in\mathcal{O}(\mathcal{V})$.\\
Since $q(\tau\mathbf{u})=q(\mathbf{u})$, fix a symmetry $\sigma$ and $\varepsilon\in\{1,-1\}$ satisfying $\sigma\tau\mathbf{u}=\varepsilon\mathbf{u}$ by Proposition \ref{prop:symmetry-matches-equal-length-nonisotropic-vectors}(iii).\\
Put $\rho=\sigma$ if $\varepsilon=1$, and $\rho=\sigma_{\mathbf{u}}\sigma$ if $\varepsilon=-1$. Then $\rho\tau\mathbf{u}=\mathbf{u}$ by Proposition \ref{prop:symmetry-matches-equal-length-nonisotropic-vectors}(i).
\item \hlt{Complement Reduction}: Put $\mathcal{W}=(\mathbb{F}\mathbf{u})^{\perp}$. For $\mathbf{x}\in\mathcal{W}$, $\innerprod{\rho\tau\mathbf{x},\mathbf{u}}=\innerprod{\rho\tau\mathbf{x},\rho\tau\mathbf{u}}=\innerprod{\mathbf{x},\mathbf{u}}=0$.\\
Hence $\rho\tau\mathcal{W}\subseteq\mathcal{W}$, and equality holds by injectivity and finite dimensionality.\\
$\mathcal{V}=\mathbb{F}\mathbf{u}\odot\mathcal{W}$ and $\dim\mathcal{W}=n-1$ by Theorem \ref{thm:metric-orthogonal-complement-dimension}(i),(ii). The space $\mathcal{W}$ is nonsingular by Theorem \ref{thm:metric-orthogonal-complement-dimension}(iv).
\item \hlt{Factor Lifting}: By induction, $(\rho\tau)|_{\mathcal{W}}$ is a product of symmetries of $\mathcal{W}$.\\
Each factor with normal $\mathbf{w}\in\mathcal{W}$ extends to $\sigma_{\mathbf{w}}\in\mathcal{O}(\mathcal{V})$, which fixes $\mathbb{F}\mathbf{u}$ and has the required restriction.\\
The lifted product equals $\rho\tau$ on $\mathbb{F}\mathbf{u}\odot\mathcal{W}$, so $\rho\tau$ is a product of symmetries.\\
The map $\rho^{-1}$ is also such a product by Proposition \ref{prop:symmetry-matches-equal-length-nonisotropic-vectors}(ii). Therefore $\tau=\rho^{-1}(\rho\tau)$ is a product of symmetries.
\end{enumerate}
\hlt{Structural Commentary}: Symmetries align one nonisotropic line; induction factors the induced complement isometry.\\
\hlt{Key Insight}: A fixed nonisotropic line exposes a lower-dimensional nonsingular complement.
\end{proof}
\end{theorem}

\begin{proposition} \label{prop:isometry-fixed-space} \hlt{Fixed Spaces of Isometries}\\
Let $\mathcal{V}$ be nonsingular and let $\tau\in\operatorname{Iso}(\mathcal{V})$. Then $\ker(I-\tau)=\operatorname{im}(I-\tau)^{\perp}$.
\begin{proof}
If $\mathbf{x}\in\ker(I-\tau)$, then $(\forall\,\mathbf{y}\in\mathcal{V})\,\innerprod{\mathbf{x},(I-\tau)\mathbf{y}}=\innerprod{\mathbf{x},\mathbf{y}}-\innerprod{\tau\mathbf{x},\tau\mathbf{y}}=0$. Hence $\ker(I-\tau)\subseteq\operatorname{im}(I-\tau)^{\perp}$.\\
$\dim(\operatorname{im}(I-\tau)^{\perp})=\dim\mathcal{V}-\dim(\operatorname{im}(I-\tau))=\dim(\ker(I-\tau))$ by Theorem \ref{thm:metric-orthogonal-complement-dimension}(i) and Theorem \ref{thm:rank-nullity}(ii). Thus $\ker(I-\tau)=\operatorname{im}(I-\tau)^{\perp}$.\\
\hlt{Structural Commentary}: Form preservation makes fixed vectors orthogonal to displacements; dimension equality closes the inclusion.\\
\hlt{Key Insight}: Fixed directions are exactly those orthogonal to every displacement.
\end{proof}
\end{proposition}

\begin{proposition} \label{prop:symmetries-fixed-hyperplanes} \hlt{Symmetries and Fixed Hyperplanes}\\
Let $\mathcal{V}$ be nonsingular orthogonal with $\operatorname{char}(\mathbb{F})\neq2$, put $q(\mathbf{x})\defeq\innerprod{\mathbf{x},\mathbf{x}}$, and let $\tau\in\mathcal{O}(\mathcal{V})$ with $\tau\neq I$.
\begin{enumerate}[label=(\roman*)]
\setlength{\itemsep}{0pt}
\item If $q(\mathbf{u})\neq0$, then $\ker(I-\sigma_{\mathbf{u}})=(\mathbb{F}\mathbf{u})^{\perp}$ is a nonsingular hyperplane.
\item The condition $(\exists\,\mathbf{u}\in\mathcal{V})\,(q(\mathbf{u})\neq0\wedge\tau=\sigma_{\mathbf{u}})$ holds if and only if $\ker(I-\tau)$ is a nonsingular hyperplane.
\end{enumerate}
\begin{proof}
\begin{enumerate}[label=(\roman*)]
\setlength{\itemsep}{0pt}
\item[$(i)$] The map $\sigma_{\mathbf{u}}$ fixes $(\mathbb{F}\mathbf{u})^{\perp}$ pointwise and sends $\mathbf{u}$ to $-\mathbf{u}$ by Proposition \ref{prop:symmetry-matches-equal-length-nonisotropic-vectors}(i).\\
The decomposition $\mathcal{V}=\mathbb{F}\mathbf{u}\odot(\mathbb{F}\mathbf{u})^{\perp}$ holds by Theorem \ref{thm:metric-orthogonal-complement-dimension}(ii). Thus $\ker(I-\sigma_{\mathbf{u}})=(\mathbb{F}\mathbf{u})^{\perp}$ by $\operatorname{char}(\mathbb{F})\neq2$.\\
This complement is a nonsingular hyperplane by Theorem \ref{thm:metric-orthogonal-complement-dimension}(i),(iv).
\item[$(ii)$] ($\Rightarrow$) If $\tau=\sigma_{\mathbf{u}}$ with $q(\mathbf{u})\neq0$, then $\ker(I-\tau)$ is a nonsingular hyperplane by (i).
\item[$(ii)$] ($\Leftarrow$) Suppose $\mathcal{H}\defeq\ker(I-\tau)$ is a nonsingular hyperplane, and put $\mathcal{L}=\mathcal{H}^{\perp}$.\\
Then $\mathcal{V}=\mathcal{H}\odot\mathcal{L}$ and $\dim\mathcal{L}=1$ by Theorem \ref{thm:metric-orthogonal-complement-dimension}(i),(ii). The line $\mathcal{L}$ is nonsingular by Theorem \ref{thm:metric-orthogonal-complement-dimension}(iv).\\
For $\mathbf{x}\in\mathcal{L}$ and $\mathbf{h}\in\mathcal{H}$, $\innerprod{\tau\mathbf{x},\mathbf{h}}=\innerprod{\tau\mathbf{x},\tau\mathbf{h}}=\innerprod{\mathbf{x},\mathbf{h}}=0$. Hence $\tau\mathcal{L}\subseteq\mathcal{L}$, and equality holds by injectivity.\\
Fix $\mathbf{0}\neq\mathbf{u}\in\mathcal{L}$ and write $\tau\mathbf{u}=\lambda\mathbf{u}$. Then $q(\mathbf{u})\neq0$ by nonsingularity of $\mathcal{L}$, while $q(\tau\mathbf{u})=q(\mathbf{u})$ implies $\lambda^2=1$.\\
The alternative $\lambda=1$ would imply $\tau=I$ on $\mathcal{H}\odot\mathcal{L}$, so $\lambda=-1$ by $\operatorname{char}(\mathbb{F})\neq2$. Finally, $\mathcal{H}=(\mathbb{F}\mathbf{u})^{\perp}$ by Theorem \ref{thm:metric-orthogonal-complement-dimension}(iii).\\
The map $\sigma_{\mathbf{u}}$ acts as $I$ on $\mathcal{H}$ and as $-I$ on $\mathcal{L}$ by Proposition \ref{prop:symmetry-matches-equal-length-nonisotropic-vectors}(i), as does $\tau$. Therefore $\tau=\sigma_{\mathbf{u}}$ by $\mathcal{V}=\mathcal{H}\odot\mathcal{L}$.
\end{enumerate}
\hlt{Structural Commentary}: Item (i) reads the fixed space; item (ii) reconstructs the flipped orthogonal line.\\
\hlt{Key Insight}: A symmetry is determined by its fixed hyperplane and normal line.
\end{proof}
\end{proposition}

\begin{proposition} \label{prop:anisotropic-space-has-no-nontrivial-unipotent-isometries} \hlt{No Unipotent Isometries in Anisotropic Spaces}\\
Let $\mathcal{V}$ be nonsingular anisotropic and let $\tau\in\operatorname{Iso}(\mathcal{V})$. If $(\exists\,m\in\mathbb{Z}_{>0})\,(\tau-I)^m=0$, then $\tau=I$.
\begin{proof}
\begin{enumerate}[label=(\roman*)]
\setlength{\itemsep}{0pt}
\item \hlt{Top Displacement}: Put $N=\tau-I$ and suppose $N\neq0$. By nilpotence, fix maximal $k\geq1$ with $N^k\neq0$.\\
Then $N^{k+1}=0$; choose $\mathbf{x}\in\mathcal{V}$ with $\mathbf{u}\defeq N^k\mathbf{x}\neq\mathbf{0}$.
\item \hlt{Adjoint Expansion}: $N^*=\tau^{-1}-I$ by the adjoint identity and Proposition \ref{prop:metric-adjoints-and-isometries}(iv).\\
Since $\tau=I+N$ and $N^{k+1}=0$, $N^*=-N+N^2-\cdots+(-1)^kN^k$ by direct multiplication.
\item \hlt{Polynomial Annihilation}: Put $P(t)=-1+t-\cdots+(-1)^kt^{k-1}$, so $N^*=NP(N)$.\\
Hence $(N^*)^kN^k=N^{2k}P(N)^k=0$ by $2k\geq k+1$ and $N^{k+1}=0$.
\item \hlt{Contradiction}: By $k$ iterations of the adjoint identity, $\innerprod{\mathbf{u},\mathbf{u}}=\innerprod{N^k\mathbf{x},N^k\mathbf{x}}=\innerprod{\mathbf{x},(N^*)^kN^k\mathbf{x}}=0$.\\
This contradicts anisotropy because $\mathbf{u}\neq\mathbf{0}$. Therefore $N=0$ and $\tau=I$.
\end{enumerate}
\hlt{Structural Commentary}: The adjoint polynomial annihilates the top displacement, forcing its self-pairing to vanish.\\
\hlt{Key Insight}: The terminal nilpotent image is forced to be isotropic.
\end{proof}
\end{proposition}

Ambient extension and orthogonal cancellation remain available by Theorem \ref{thm:witt-extension-theorem-metric-spaces} and Corollary \ref{cor:witt-cancellation-metric-spaces}.
